\section{Introduction}
One of the oldest and most prominent economic narratives holds that swings in investor beliefs drive asset prices and business cycles. From \citet{sprague1910history} and \citet{fisher1933debt}, to \citet{keynes1936general}, \citet*{minsky1986stabilizing}, \citet*{kindleberger1996history}, and \citet*{Sh00}, a common theme emerges: fueled by leverage, a wave of excessive optimism drives up stock prices, spurring excessive firm investment and hiring. Once earnings disappoint, asset prices collapse on Wall Street with a blast wave felt all the way onto Main Street. Despite the prominence of this narrative, it remains unclear what ingredients theory needs to formalize these narratives in a quantitatively meaningful way.

We formalize these narratives in an otherwise frictionless economy. Markets are complete, prices are flexible, and there are no financial constraints. The only link between investor beliefs and output is that labor hiring is inherently risky: because employment must be chosen before productivity is realized, labor demand becomes a risky investment. The firm's discount factor---a wealth-weighted average of investor beliefs---dictates how this hiring risk is priced. When optimists grow wealthier, risk premia are compressed, and the firm discounts risk less aggressively, boosting labor demand. When this optimism is overcome by pessimism, the opposite occurs. Thus, waves of optimism on Wall Street indeed percolate onto Main Street.

Beyond formalizing the belief-driven business cycle narrative, the model yields sharp theoretical predictions. We derive an analytic taxonomy of belief systems that allows us to classify these systems according to their business-cycle amplification, risk build-up, and turnover properties. Moreover, disciplined by survey data on earnings expectations, belief heterogeneity accounts for these narratives in a quantitatively meaningful way. The model jointly matches:
\begin{enumerate}
    \item \textit{Asset-pricing and business-cycle moments.} The equity premium, stock-return volatility, and the volatility of hours worked---moments that are typically addressed in isolation but share a common origin in our setting: time-varying risk premia driven by shifts in the wealth distribution across investors.

    \item \textit{Excess volatility and return predictability.} Returns are more volatile than dividends, and the price-dividend ratio predicts future returns with a Campbell--Shiller discount-rate share close to unity, as in the data.

    \item \textit{Turnover dynamics.} Belief dispersion drives trading volume, and this relationship strengthens substantially in recessions---particularly around state transitions---consistent with patterns we document using I/B/E/S analyst forecasts and CRSP turnover data.

    \item \textit{Forecast biases in productivity expectations.} The model generates over-extrapolation in productivity growth forecasts, with a Coibion--Gorodnichenko coefficient that closely matches our estimates from the Survey of Professional Forecasters.
\end{enumerate}

\paragraph{The Model.} 
In the model, households differ in their beliefs about the evolution of TFP. Given their distinct views, they actively trade stocks, affecting asset prices. They supply labor with preferences that eliminate wealth effects, as in \citet[henceforth GHH]{GHH88}. 

The key feature is that heterogeneity in beliefs shapes the firm's labor demand because hiring is risky. As in \citet{burnside1993labor}, labor is chosen before the shock realizations. This timing induces an operating leverage channel whereby hiring becomes a risky investment, as it can add cost pressure to firms, which also occurs in labor search models.\footnote{We formally show that this form of \emph{investment in labor} is a simplified version of the investment in hours that happens in labor-search models.} This feature is critical to linking business-cycle fluctuations to asset-price fluctuations. We focus on the firm's investment in forming a workforce, as opposed to physical capital formation, because it is understood that, absent additional frictions, fluctuations in capital investment cannot be a source of business cycles; instead, what is needed are movements in the labor wedge \citep{CKM07}.\footnote{Investment is small relative to the capital stock in a business-cycle model. Thus, fluctuations in investment do not meaningfully impact the production possibility frontier.} Moreover, the operating leverage channel leads to testable implications regarding the correlation of returns and the labor share that is consistent with data, \citep[e.g., see][]{donangelo2019cross}.

Belief heterogeneity is essential for generating leverage and turnover dynamics, both recurrent themes in optimism-wave narratives. Because of heterogeneous beliefs, households take long or short positions in the firm. Optimistic households increase leverage, amplifying their risk exposure. Leverage creates internal propagation by magnifying differences in investors' risk exposures: optimistic investors accumulate relatively more wealth when positive shocks occur, while pessimists gain during downturns.

These dynamics produce internal propagation as wealth-weighted beliefs jointly affect asset prices and output through the operating leverage channel. When wealth-weighted beliefs become more optimistic, increased demand for risky assets pushes stock prices upward, compresses risk premia, and encourages hiring. Thus, as investors' risk appetite increases, discount rates decline, prompting more hiring by a firm attempting to raise its market value. Through these dynamics, the \emph{equity premium}, the \emph{equity volatility}, and \emph{labor volatility} puzzles are tied together into a single volatility puzzle. 

\paragraph{Theoretical Results.} Despite featuring Epstein-Zin preferences and endogenous labor supply, the model remains highly tractable due to a convenient ``as if'' property portable to other settings. Specifically, the equilibrium is characterized as if the economy were an endowment economy, with a common stochastic discount factor (SDF) used to price assets. Importantly, the firm makes risky hiring decisions using a common SDF. To increase the firm's value, it constructs a risk-neutral productivity growth expectation consistent with the SDF. This risk-neutral expectation ultimately determines employment. Crucially, we impose no assumptions regarding the firm's direct knowledge of investor beliefs; the firm only needs to understand how its hiring decisions affect its market value. 

In the special case of log utility, the model admits an analytic solution, enabling us to articulate a clear taxonomy of belief systems. \q{ResultSummary}{Specifically, we classify beliefs into two broad categories, each with two sub-categories. The first category comprises \emph{rank-preserving} beliefs, subdivided into optimistic and pessimistic beliefs. The second category includes \emph{rank-alternating} beliefs, subdivided into extrapolative and intrapolative beliefs. This taxonomy allows us to derive general properties regarding how each belief structure affects business-cycle dynamics.

We identify several principles regarding the amplification properties under different belief systems. First, we show that only extrapolative beliefs amplify business-cycle fluctuations across all states of the economy relative to rational expectations. Amplification occurs because extrapolation boosts the firm's discounting of risk states during booms but depresses it during busts, a property unique to extrapolative beliefs.

Second, we uncover that only extrapolative beliefs engender \textit{risk build-up}, i.e., longer booms lead to deeper busts. This phenomenon, also exclusive to extrapolative belief systems, occurs because extrapolative households consistently accumulate wealth throughout boom phases. When the economy transitions into a downturn, extrapolative agents remain wealthier while simultaneously becoming the most pessimistic investors. The reversal in the wealthiest households' optimism significantly amplifies the economy's rift.}

Third, we demonstrate that belief heterogeneity drives stock turnover through two first-order forces: by inducing passive rebalancing as differences in beliefs translate into differences in wealth growth rates within states and by the changes in beliefs when the economy changes states. We further show that when beliefs are rank-alternating, turnover increases most when the economy enters a recession, a property that we verify in the data.

Our taxonomy highlights the importance of belief heterogeneity for understanding macro-finance dynamics. Empirical evidence from \citet{lopez2017credit} and \citet{krishnamurthy2017credit} indicates that risk premia are low during credit booms preceding severe crashes.\footnote{See also \citet{gennaioli2018crisis} and \citet{krishnamurthy2025dissecting}.
} In contrast, standard macro-finance models without belief heterogeneity predict precisely the opposite  \citep[e.g.,][]{brunnermeier2014macroeconomic,he2013intermediary}.
Furthermore, our analysis helps distinguish models driven by heterogeneous beliefs from those driven by heterogeneous risk aversion \citep[e.g.,][]{panageas2020implications}, a challenge for the literature. Our results clarify that heterogeneity in risk aversion cannot generate risk build-up germane to extrapolative belief systems. 

\paragraph{Quantitative Results.}
To complement our theoretical insights, we provide a quantitative evaluation. Aside from beliefs, our framework follows canonical assumptions; we adopt a standard calibration approach for preferences and the process for TFP. To discipline the calibration of beliefs, we utilize survey data on earnings forecasts. The belief process is parsimonious: two key parameters---an overreaction coefficient and the volatility of a sentiment shock---are pinned down by the degree of return extrapolation observed in surveys and the excess volatility of survey-based expectations relative to realized outcomes. This specification nests the belief systems identified in our taxonomy, so we let the data determine whether the empirically relevant system is extrapolative, rank-preserving, or a combination.

Our quantitative analysis underscores the importance of each ingredient by following a progression. Operating leverage alone matches dividend-growth volatility but fails on equity returns and employment. Homogeneous subjective beliefs add labor fluctuations but cannot generate a realistic equity premium. Only the interaction of rational and extrapolative investors---and the wealth dynamics it produces---allows the model to jointly match the moments listed above.

\paragraph{Literature Review.}

This paper connects to the literature initiated by \citet{KP82} and \citet{MP85}, who used complete-market benchmarks to analyze business cycle fluctuations and asset prices in settings where these two dimensions do not feed into each other. A recent strand emphasizes the role of fluctuating risk premia in driving real cycles, particularly highlighting the riskiness of labor hiring \citep[e.g.,][]{H17, BB19, DH19, KLMP19}. We bridge this channel with the literature on heterogeneous beliefs and asset pricing.

\citet{HK78} and \citet{SX03} formalized how heterogeneous beliefs induce asset-market trade, emphasizing that belief heterogeneity engenders asset-price volatility through leverage. A virtue of these formalizations is their subtle departure from strict rational expectations: agents fully understand the structure of the economy and the risks it entails, yet disagree about the distribution of future fundamentals. Investors can therefore hold different views about which firms or technologies are more likely to succeed, while recognizing the risks inherent in taking positions on those outcomes. Importantly, a modeler is not free to impose arbitrary beliefs about asset prices onto agents whose decisions determine those very prices. We embrace this approach. 

Building on those foundations, a large body of work examines how belief heterogeneity interacts with financial constraints to generate speculative bubbles \citep[e.g.,][]{G03, FG08, G10, S13, INS18, barlevy2014leverage, barlevy2022speculative}, while closely related papers explore how wealth redistribution driven by heterogeneous beliefs affects asset-price dynamics \citep{DM94, XY09, KS12, wr2022sentiment}. This earlier literature set aside the connection to the real economy. We view our paper as a natural progression of these models, linking belief heterogeneity to the real economy through the labor market.

A growing body of macro-finance work embeds beliefs into environments where these affect production, but optimism affects output by interacting with market imperfections such as incomplete markets, financial frictions, or demand externalities. Some models connect speculative bubbles to real activity through capital investment,\footnote{See \cite{GHH05}, \cite{BSX06}, \cite{P05}, and \cite{BDUV16}.} while others explore how heterogeneous beliefs interact with aggregate demand externalities, highlighting how changes in wealth distribution can amplify demand-driven recessions \citep{CS19, CS20, guerreiro2022disagreement, CS24}. \citet{simsek2021macroeconomics} provides a comprehensive review.

Our paper fills a gap between these two literatures: models where belief heterogeneity drives asset prices but has no impact on the real economy, and models where beliefs affect output only by aggravating market imperfections. To our knowledge, there is no frictionless benchmark economy where belief heterogeneity is a source of business-cycle fluctuations in its own right. Developing such a benchmark is important: First, by abstracting away market imperfections, we can provide a qualitative and quantitative assessment of the direct effects of beliefs on the business cycle. Second, we can uncover general principles that carry over to environments with market imperfections. Finally, a frictionless benchmark serves as an ideal target that policy interventions to mitigate market imperfections should strive to replicate. After all, policy should not care about asset prices per se, but only inasmuch as these create inefficiencies.

Our theoretical results emphasize that extrapolative beliefs conform best with the narratives and empirical patterns documented above, providing a disciplined rationale for their prominence in the literature. This connects us to models featuring diagnostic expectations \citep{GS10, bianchi2024diagnostic} and, more broadly, to the literature on subjective beliefs and macroeconomic dynamics \citep[e.g.,][]{EP11, ACD18, BGS18, BBH19}. Most closely related, \cite{AM19} shows how extrapolative beliefs simultaneously explain stock market and business cycle fluctuations, but in a setting with homogeneous beliefs. Our taxonomy clarifies that heterogeneity is essential for generating risk build-up and the turnover dynamics that are central to optimism-wave narratives.

Finally, our model generates sharp predictions for trading volume, a dimension that receives surprisingly little attention in equilibrium asset-pricing models.\footnote{\citet{AlvarezAtkeson2018} develop a general equilibrium model where trade volume becomes a pricing factor through idiosyncratic shocks to risk tolerance.} On the empirical side, \citet{vanderBeck2025} show that the ratio of return volatility to portfolio turnover provides a lower bound on price impact that is tightened by disagreement, establishing that trading activity is central to understanding asset price movements. \citet{ThesmarVerner2025} and \citet{McCarthyHillenbrand2025} independently document that disagreement between sophisticated and extrapolative investors is associated with higher trading volume and accounts for a large share of stock price variation. Our model complements these findings with a structural foundation for these findings, linking disagreement to turnover, asset prices, and employment within a unified framework.

\paragraph{Roadmap.} The next section presents three empirical facts about investor beliefs documented in the literature to motivate our analysis. Subsequently, we present the theoretical model followed by the quantitative evaluation. Finally, we conclude.